\section{Diagnósticos de ajuste: gráficos QQ e resíduos quantílicos}
\label{ap:diagnosticos}

\setcounter{figure}{0}
\renewcommand{\thefigure}{A.\arabic{figure}}

Este apêndice reúne os gráficos de diagnóstico técnico do documento,
para que o corpo do texto flua sem interrupções; cada figura mantém o
seu bloco de leitura. Duas noções bastam para acompanhar tudo.
\textbf{Gráfico QQ}: compara os dados reais com um molde matemático de
distribuição; se os pontos seguem a reta, o molde descreve bem os
dados, e desvios sistemáticos da reta indicam onde o molde falha. O
número r que acompanha cada QQ resume essa aderência: quanto mais
próximo de 1, melhor. \textbf{Resíduo quantílico}: é a versão do QQ
para um modelo já ajustado com covariáveis; se o modelo descreve bem
os dados, esses resíduos se comportam como uma Normal padrão e os
pontos deitam na diagonal.

\subsection*{Diagnósticos da fase exploratória}

As Figuras~\ref{fig:qqcusto} e~\ref{fig:qqtmp} avaliam os moldes
marginais de faturamento e TMP da Seção~\ref{sec:distribuicoes}.

\begin{figure}[H]
\centering
\subfig{fig_ae_03_qq_lognormal_custo_saida.png}{Molde LogNormal}
\subfig{fig_ae_03_qq_fisk_custo_saida.png}{Molde log logístico (Fisk)}
\caption{Gráficos QQ do faturamento por saída: os pontos são os dados e a
reta é o comportamento esperado pelo molde.}
\label{fig:qqcusto}
\end{figure}

\comoler{quanto mais os pontos abraçam a reta, melhor o molde. No
painel (b), o molde log logístico segue os dados por quase toda a
extensão, descolando apenas nos pouquíssimos hospitais mais caros; por
isso ele recebe a melhor nota da Tabela~\ref{tab:familias}.}

\begin{figure}[H]
\centering
\subfig{fig_ae_03_qq_lognormal_tmp.png}{Molde LogNormal}
\subfig{fig_ae_03_qq_fisk_tmp.png}{Molde log logístico (Fisk)}
\caption{Gráficos QQ do TMP no painel completo.}
\label{fig:qqtmp}
\end{figure}

\comoler{nos dois painéis os pontos acompanham a reta em quase toda a
extensão e descolam apenas nos pouquíssimos hospitais de permanência
mais longa. O painel corrente já não contém a subpopulação de longa
permanência que, na versão anterior desta análise, criava um teto de
30,5 dias e o característico achatamento no alto: a ETAPA F a removeu, e
o ajuste do TMP passou a ser quase perfeito.}

As Figuras~\ref{fig:qqcondA} a~\ref{fig:qqcondC} sustentam a conclusão
da Seção~\ref{sec:distribuicoes} de que o desajuste dos QQ marginais
vem de misturar 289 hospitais heterogêneos, não do molde.

\begin{figure}[H]
\centering
\subfig{fig_ae_11_qq_custo_marginal.png}{Faturamento: todos os hospitais
juntos}
\subfig{fig_ae_11_qq_custo_residuo_ef.png}{Faturamento: cada hospital
comparado consigo mesmo}
\caption{Faturamento por saída: gráfico QQ antes e depois de controlar as
diferenças fixas entre hospitais e entre anos.}
\label{fig:qqcondA}
\end{figure}

\comoler{compare o gráfico de cima com o de baixo: em cima, todos os
hospitais juntos, os pontos serpenteiam em volta da reta; embaixo,
comparando cada hospital consigo mesmo, os pontos praticamente deitam
na reta.}

\begin{figure}[H]
\centering
\subfig{fig_ae_11_qq_tmp_marginal.png}{TMP: todos os hospitais juntos}
\subfig{fig_ae_11_qq_tmp_marginal_sem_longa_perm.png}{TMP: sem a antiga
subpopulação de longa permanência}
\caption{TMP: gráfico QQ do painel completo e da versão sem a antiga
subpopulação de longa permanência --- no painel corrente as duas
coincidem, porque a ETAPA F já removeu essa subpopulação.}
\label{fig:qqcondB}
\end{figure}

\comoler{no painel anterior este par contrastava o desenho em degrau
(com os hospitais de longa permanência) com a reta quase perfeita (sem
eles). No painel corrente a ETAPA F já removeu essa subpopulação, então
as duas versões coincidem e ambas mostram a reta --- a confirmação de
que o desvio era a subpopulação, não o molde.}

\begin{figure}[H]
\centering
\subfig{fig_ae_11_qq_mort_logito_marginal.png}{Mortalidade: todos os
hospitais juntos}
\subfig{fig_ae_11_qq_mort_logito_residuo_ef.png}{Mortalidade: cada
hospital comparado consigo mesmo}
\caption{Mortalidade: gráfico QQ antes e depois de controlar as
diferenças fixas entre hospitais e entre anos.}
\label{fig:qqcondC}
\end{figure}

\comoler{mesmo padrão do faturamento: a curva em S do gráfico de cima, com
todos os hospitais misturados, vira uma reta bem comportada quando cada
hospital é comparado consigo mesmo. A heterogeneidade entre hospitais,
e não o molde, era a fonte do desajuste.}

\subsection*{Resíduos quantílicos dos modelos estimados}

As Figuras~\ref{fig:qqprimeira} e~\ref{fig:qqsegunda} validam as
famílias dos modelos da Seção~\ref{sec:resultados}, e a
Figura~\ref{fig:qqgama} documenta a reestimação da UTI comentada na
Seção~\ref{sec:inferenciafinal}.

\begin{figure}[H]
\centering
\subfig{fig_est_04_qq_mort_zib.png}{Mortalidade geral, modelo Beta com
inflação em zero}
\subfig{fig_est_04_qq_pct_alta_zoib.png}{Fração de alta complexidade,
modelo ZOIB}
\caption{Resíduos quantílicos dos modelos de taxa.}
\label{fig:qqprimeira}
\end{figure}

\begin{figure}[H]
\centering
\subfig{fig_est_04_qq_producao_bn.png}{Produção, Binomial Negativa}
\subfig{fig_est_04_qq_ocup_uti_pos.png}{Ocupação de UTI entre os
positivos, o único ajuste insatisfatório}
\caption{Resíduos quantílicos da produção e da parte positiva da UTI.}
\label{fig:qqsegunda}
\end{figure}

\comoler{quanto mais os pontos azuis acompanham a linha tracejada,
melhor o modelo descreve os dados; o valor de r no canto é essa
aderência resumida em um número (1 seria perfeito). Compare o
alinhamento quase exato da fração de alta complexidade com o
descolamento da cauda esquerda na UTI, que motivou a reestimação da
figura seguinte.}

\begin{figure}[H]
\centering
\includegraphics[width=.72\textwidth]{fig_est_08_qq_uti_gama.png}
\caption{Resíduos quantílicos da intensidade de UTI reestimada com a
família Gama: a aderência sobe de 0,814 para 0,944.}
\label{fig:qqgama}
\end{figure}

\comoler{compare com o painel (b) da figura anterior: os pontos que
antes despencavam da linha tracejada na ponta esquerda agora a
acompanham em quase todo o trajeto.}
